\subsection{C.2 supplement. A computed coboundary with a different boundary}

\begin{proposition}[C.2.2 (Proved): the untwisted intersection form]
Let $W_A$ be a sufficiently small compact tubular representative of the
smooth total space $\mathcal X^+$ of A.2.2 over a closed $u$-disc.
Its boundary is the mapping torus of Engel's \emph{untwisted} family
$X(u)$. It is not the six-sphere mapping torus $M_\Phi$.
Write $\mathcal L=H^2(S;\Z)$, $F=c_1\mathcal O_S(S_\infty)$,
$m=c_1(L)$, and $v=6P-O-3F$ in the conventions of A.4.1.
Let $H$ be the infinity-section class on $R$ and $A$ the exceptional
Cartier divisor pulled back to $R^+$. Then
\[
 H_4(W_A;\Z)=
 \frac{\mathcal L\oplus\Z H\oplus\Z A}{\Z(m-A)}.
\]
For $\alpha=(a,h,b)\in H^2(R^+;\Z)$ set
\[
 d(\alpha)=(I-T^*)a+h(m-F)+bF.
\]
The middle intersection form of $W_A$ is
\[
 B(\alpha,\beta)=-d(\alpha)\cdot d(\beta).
\]
Its rank is two, its signature is $+2$, and its two nonzero Smith
factors are $1,11$. On the image modulo its radical it is represented,
in the ordered basis $(v,m)$, by
\[
 \begin{pmatrix}6&1\\1&2\end{pmatrix}.
\]
\end{proposition}
\begin{proof}
The local semistable charts $xy=u$ give a tubular representative which
retracts to $Q^+$; its boundary is the nearby-fibre mapping torus. One
may equivalently choose a regular neighbourhood of $Q^+$ in the smooth
total space and round its boundary. The fibrewise plumbing identifies
that boundary with the inverse image of a small positive $u$-circle.
The normalization mapping cylinder and $H_3(S)=0$ identify its fourth
homology with the cokernel of
$H_4(S)\to H_4(R^+)$. In the divisor coordinates of A.4.2 the two
sections are $D'_0=H+m-A$ and $D_\infty=H$. Their difference is
$m-A$, proving the group presentation.

The normalization is an immersed complex hypersurface in $W_A$.
Its normal line has class $-D'_0-D_\infty$, because the whole central
divisor is the divisor of the base coordinate $u$. Put
$r_0\alpha=a+bF$ and $r_\infty\alpha=a-hm$; these are the
restrictions to the two sections. The self-intersection of two
four-cycles pushed forward from the normalization contains the normal
Euler term and the intersections between distinct local branches.
All transverse complex branch intersections have positive sign.
Consequently it is
\begin{align*}
 &-r_0\alpha\cdot r_0\beta-r_\infty\alpha\cdot r_\infty\beta\\
 &\quad+r_0\alpha\cdot T^*r_\infty\beta
       +T^*r_\infty\alpha\cdot r_0\beta\\
 &=-\bigl(r_0\alpha-T^*r_\infty\alpha\bigr)
       \cdot\bigl(r_0\beta-T^*r_\infty\beta\bigr).
\end{align*}
This is an intersection-of-cycles calculation in the actual smooth
fourfold, rather than an assignment of a form to an unconstructed
central fibre. The isometry property of $T$ and $T^*m=m-F$ give
the stated expression for $d$. Also $d(m-A)=0$, so it descends.

The translation formula in A.4.1 shows that the image of $d$ is
$\langle F,v,m\rangle$. The surjectivity of the three intersection
coordinates proved there shows equality integrally. Its radical is
$\Z F$, since $F^2=Fv=Fm=0$ and
$v^2=-6$, $m^2=-2$, $vm=-1$. The remaining form is the displayed
positive matrix, of determinant eleven and with an entry equal to
one. Integral row and column reduction gives Smith factors $1,11$.
These computations, including the full eleven-by-eleven form, are
independently checked in \path{code/e_push_C_attempts.py}.
\end{proof}

\begin{proposition}[C.2.3 (Proved): a relative Pontryagin number for $W_A$]
In the coordinates of C.2.2, $p=p_1(TW_A)$ evaluates as
\[
 p(a,h,b)=-4(aF)-28h.
\]
It vanishes on the radical of the middle form and therefore has a
rational relative lift. The square of any such lift is
\[
 \bar p^{\,2}=\frac{4512}{11}.
\]
These numbers do not compute $\mu(M_\Phi)$. In particular the
canonical determinant class of $W_A$ restricts nontrivially to its
own boundary, so C.2.1 cannot be applied with that determinant.
\end{proposition}
\begin{proof}
Let $h:R^+\to R$ be the modification, $C$ its nodal elliptic-fibre
centre, and $c=c_1(R)=F+2H+m$. The centre is a regular embedding of
codimension two, also at its node: its ideal there is $(X,zv)$.
Its normal bundle has degree zero, and $c\cdot C=0$. The blow-up
and small resolution give
\[
 c_1(R^+)=h^*c-A,\qquad
 A\,h^*D\,h^*E=0,\quad
 A^2h^*D=-C\cdot D,\quad A^3=0.
\]
These are the regular-embedding blow-up identities. The subsequent
small resolution is crepant, and $A$ is the pullback of a Cartier
divisor, so they remain valid on $R^+$. The normal line to the
immersed normalization has class
$h^*N_0+A$, where $N_0=-(2H+m)$.

For completeness, the required second Chern pairings do not assume
that $R^+$ is K\"ahler. Analytic Hirzebruch--Riemann--Roch applies to
a compact complex manifold; see \citet[Theorem~4.3, pp.~564--565]{AtiyahSinger1968IIIE}.
The modification has $Rh_*\mathcal O_{R^+}=\mathcal O_R$.
For the blow-up this follows locally from its hypersurface description
in the projective line bundle and the vanishing of the direct images
of $\mathcal O_{\mathbb P^1}(-1)$; for the small resolution it follows
from the formal neighbourhood of its $(-1,-1)$ rational curve.
Thus $\chi(h^*D)=\chi(D)$ for every line bundle $D$ on $R$.
The linear coefficient in the Riemann--Roch polynomial for $kD$
and the preceding intersection identities give
\[
 c_2(R^+)h^*D=c_2(R)D+C\cdot D.
\]
Moreover $\chi(\mathcal O(A))=\chi(\mathcal O)=1$: on the
exceptional projective line bundle over $C$, the restriction is its
relative $\mathcal O(-1)$, which has zero derived direct image.
Pulling through the small resolution does not change this Euler
characteristic. Applying Riemann--Roch to $A$ yields $c_2(R^+)A=0$,
because $A^3$, $c_1(R^+)A^2$, and $c_1(R^+)^2A$ all vanish.

The real tangent bundle of $W_A$ restricted to the normalization is
the sum of its real tangent bundle and this normal line. Therefore
\[
 h^*p_1(R)+h^*N_0^2-4[C]
\]
has the same pairings on pulled-back divisors as $p_1(TW_A)$, and
the pairing with $A$ is zero. Here $[C]$ denotes its degree-four
cohomology class on $R$. Since $H^2=-mH$, we have
$(2H+m)^2=m^2=-2[\mathrm{pt}]$. The surface $S$ has
$p_1(S)=-24[\mathrm{pt}]$. The ruled-bundle tangent sequence thus
gives
\[
 p_1(R)+N_0^2=p_1(S)+2(2H+m)^2=-28[\mathrm{pt}],
\]
and $C\cdot(a+hH)=aF$. This proves the covector formula.

Modulo $\Z F$ the coordinate of $d(a,h,b)$ along $v$ is $aF$
and the coordinate along $m$ is $h$. Thus the covector on the
nondegenerate quotient is $(-4,-28)$. It annihilates the radical,
which is precisely the condition for a rational relative lift.
Its square is
\[
 (-4,-28)\frac1{11}
 \begin{pmatrix}2&-1\\-1&6\end{pmatrix}
 \binom{-4}{-28}=\frac{4512}{11}.
\]
Different relative lifts differ by a class pairing trivially with
absolute classes, so this value does not depend on the choice.

Finally the canonical first Chern class on the nearby untwisted
fibre is $f^*c_1\mathcal O_B(1)$. It is nonzero: the surjectivity of
$a\mapsto(aF,av,am)$ in A.4.1 supplies an integral two-cycle with
evaluation one, and the Mayer--Vietoris relations in A.4.2 do not
alter that evaluation. Hence no boundary trivialization of this
determinant exists. Even the formal expression
$(\bar p^{\,2}-4\sigma(W_A))/896=79/176$ is only an arithmetic
combination for this different boundary, not the invariant requested
in Part C. A cobordism implementing the finite logarithmic transforms
and the cusp closing map would be required to transport the result.
\end{proof}
